\documentclass[10pt,twocolumn]{article}
\usepackage[scaled]{helvet}
\renewcommand\familydefault{\sfdefault}
\usepackage[T1]{fontenc}
\usepackage[margin=0.7in]{geometry}
\usepackage{titlesec}
\usepackage{enumitem}
\usepackage{parskip}
\setlength{\columnsep}{0.24in}
\setlist[itemize]{leftmargin=1.2em, topsep=0.2em, itemsep=0.18em}
\titleformat{\section}{\large\bfseries}{\thesection}{1em}{}

\begin{document}
\title{\vspace{-1.6cm}AWS ECR (Elastic Container Registry)}
\author{AWS ML Study Notes}
\date{}
\maketitle
\small

\section*{Overview}
ECR is AWS managed container image storage. It holds versioned Docker and OCI images that other services can pull when they need to start a container.

For ML teams, ECR is where custom training images, inference images, and utility containers are stored so that deployments are reproducible and centrally managed.

\section*{Core Concepts}
\begin{itemize}
\item Repositories group related images, while tags and digests identify specific builds. Digests are immutable and better for strict reproducibility than floating tags.
\item ECR supports vulnerability scanning, lifecycle policies, encryption, and repository policies for cross account sharing.
\item Authentication is integrated with IAM. Build systems and runtimes usually obtain temporary credentials and then push or pull images through the ECR API.
\end{itemize}

\section*{How It Works}
\begin{itemize}
\item A CI system builds the image, runs tests, tags it, and pushes it to an ECR repository.
\item Services such as ECS, SageMaker, or Lambda container images then pull the image when creating or updating a workload.
\item Multi account setups often replicate approved images from a build account into runtime accounts to separate software production from software consumption.
\end{itemize}

\section*{AI and ML Relevance}
\begin{itemize}
\item ML inference containers often include model servers, preprocessing logic, and optimized native libraries. ECR ensures those dependencies are packaged consistently across environments.
\item Training images in ECR let you standardize frameworks, CUDA libraries, and internal tooling for repeatable experiments and pipelines.
\end{itemize}

\section*{Operational Notes}
\begin{itemize}
\item Avoid using mutable tags such as latest as the only deployment reference. Use digests or versioned tags in release automation.
\item Enable lifecycle policies so old image layers do not accumulate forever and quietly increase storage cost.
\item Scan images, patch base layers regularly, and keep the build chain deterministic so security and debugging stay tractable.
\end{itemize}

\section*{What To Remember}
\begin{itemize}
\item ECR stores images; ECS, EKS, SageMaker, and Lambda run them.
\item Tagging is for convenience. Digests are for certainty.
\item Treat the container registry as part of the software supply chain, not just as storage.
\end{itemize}

\end{document}
